%%%%%%%%%%%%%%%%%%%%%%%%%%%%%%%%%%%%%%%%%%%%%%%%%%%%%%%%%%%%%%%%%%%%%%%%%%%
% SUCESIONES ALICUATORIAS,
% por
% MANUEL BENITO y JUAN LUIS VARONA,
% escrito para La Gaceta de la RSME, por petici\'on de J. Cilleruelo.
%     FECHA COMIENZO ESCRITURA: 20 de marzo de 1999.
%     FECHA \'ULTIMA MODIFICACI\'ON: 27 de mayo de 1999.
%     NOTA POSTERIOR: 2 de agosto de 2000.
% Escrito en LaTeX2e, por Juan Luis Varona (para problemas del texto,
% contactar con jvarona@dmc.unirioja.es).
% Formato, tama\~no del "body" de texto y espaciado seg\'un
% indicaci\'on de J. Cilleruelo (franciscojavier.cilleruelo@uam.es).
%%%%%%%%%%%%%%%%%%%%%%%%%%%%%%%%%%%%%%%%%%%%%%%%%%%%%%%%%%%%%%%%%%%%%%%%%%%

%% ESTO ES LO QUE ME DIJO CILLERUELO:
%% \documentclass[11pt,a4paper,twoside]{report}

%% DE MOMENTO, YO USO ESTO OTRO:
\documentclass[11pt,a4paper,twoside]{article}

\setlength{\textheight}{19.5cm}
\setlength{\textwidth}{13cm}

\frenchspacing

\baselineskip=12pt \par % No pinta nada, pues es lo que viene por defecto.
% Adem\'as, si se cambiara as\'{\i} desde LaTeX ser\'{\i}a una chapuza.
% ESTOS CAMBIOS HAY QUE HACERLOS DESDE UNA CLASE,
% Y QUE EL DOCUMENTO CARGUE LA CLASE.
% Adem\'as, aqu\'{\i} no tiene efecto, pues el \begin{document} carga la
% suya (la de la clase que se est\'e usando) y pasa de \'esta. Para que
% tuviera efecto, tedr\'{\i}a que estar despu\'es del \begin{document}.
% Meter formato despu\'es de \begin{document} va en contra de toda la
% filosof\'{\i}a de estructuraci\'on de LaTeX.
%\noindent\begin{tabular}{lr}
\itemsep=0pt \par
\parsep=0pt \par
% EN LAS SEPARACIONES, ES MEJOR A\~NADIR "PLUS" Y "MINUS".
% EL "	PLUS" Y EL "MINUS" HACEN DE GOMA PARA COMERSE LOS OVER/UNDERFULL
\parindent=7mm

\usepackage[spanish]{babel}

\usepackage{amssymb}
%\usepackage{amstext} % Incluida en el paquete "amsmath"
\usepackage{amsmath}

\usepackage{url}
% SE PUEDE COGER EN:
% ftp://CTAN/tex-archive/macros/latex/contrib/other/misc/url.sty
% SE USA PARA LOGRAR CORTES DE PALABRA EN PAGINAS WEB Y SIMILARES.

%%%%%%%%%%%%%%%%%%%%%%%%%%%%%%%%%%%%%%%%%%%%%%%%%%%%%%%%%%%%%%%%%%%%%%%%%%%

% YA NO USO NINGUNA DE LAS DOS EN LA VERSION REDUCIDA
% Definiciones originales de \bmod y \pmod originales en p. 361 del TeXbook.
% Aqu\'{\i} las redefinimos para que mod tenga acento en la o.
%\renewcommand{\bmod}{\mskip-\medmuskip \mkern5mu
%  \mathbin{\rm m\acute od} \penalty900 \mkern5mu \mskip-\medmuskip}
%\renewcommand{\pmod}[1]{\allowbreak \mkern18mu ({\rm m\acute od}\,\,#1)}
% Realmente, s\'olo vamos a usar \bmod.

% TAMPOCO LA USO EN LA VERSION REDUCIDA
%\newcommand{\mcd}{\mathop{\rm mcd}}

% ESTAS DOS DEFINICIONES TIENEN INFORMACION QUE NO USAMOS EN ESTE ARTICULO
% \newcommand{\info}{\(\begin{array}{@{}c@{}}{n\strut}\\{k}\\{\text{d\'{\i}gitos}}
% \\{\text{gu\'{\i}a}}\end{array}\)}
% \newcommand{\num}[4]{\(\begin{array}{@{}c@{}}{#1\strut}\\{#2}\\{#3}\\{#4}
% \end{array}\)}

% EN ESTE ARTICULO USAMOS ESTAS (omitimos la informaci\'on del gu\'{\i}a):
\newcommand{\info}{\(\begin{array}{@{}c@{}}{n\strut}\\{k}\\{\text{d\'{\i}gitos}}
\end{array}\)}
\newcommand{\num}[4]{\(\begin{array}{@{}c@{}}{#1\strut}\\{#2}\\{#3}
\end{array}\)}

\newcommand{\pp}{\!\cdot\!}

% ESTO TAMOCO LO USAMOS:
%\newcommand{\nodriver}{{}^{(*)}} % Marcar los que no est\'an en "driver"
%\newcommand{\ldw}[1]{\mbox{\scriptsize\ (#1)}} % Ultimo driver con [1] digitos
% O SEA QUE PONEMOS ESTO OTRO:
\newcommand{\nodriver}{}
\newcommand{\ldw}[1]{}

%%%%% REDEFINICIONES PARA ESPACIADO DE TABLAS FLOTANTES COMO EN "amsart"
\setcounter{topnumber}{4}
\setcounter{bottomnumber}{4}
\setcounter{totalnumber}{4}
\setcounter{dbltopnumber}{4}
\renewcommand{\topfraction}{.97}
\renewcommand{\bottomfraction}{.97}
\renewcommand{\textfraction}{.03}
\renewcommand{\floatpagefraction}{.9}
\renewcommand{\dbltopfraction}{.97}
\renewcommand{\dblfloatpagefraction}{.9}
\setlength{\floatsep}{12pt plus 6pt minus 4pt}
\setlength{\textfloatsep}{15pt plus 8pt minus 5pt}
\setlength{\intextsep}{12pt plus 6pt minus 4pt}
\setlength{\dblfloatsep}{12pt plus 6pt minus 4pt}
\setlength{\dbltextfloatsep}{15pt plus 8pt minus 5pt}
% Sacados de "amsart.cls".
% Ver el significado en p. 142 del "LaTeX Companion".

%%%%%%%%%%%%%%%%%%%%%%%%%%%%%%%%%%%%%%%%%%%%%%%%%%%%%%%%%%%%%%%%%%%%%%%%%%%

\title{Sucesiones alicuatorias
\footnote{ESTE ART\'ICULO APARECI\'O PUBLICADO EN: La Gaceta de la Real Sociedad
Matem\'atica Espa\~nola, Vol.~2, n.\textsuperscript{o}~2 (1999), 357--365.}
}

\author{Manuel Benito y Juan Luis Varona}

\date{}

%%%%%%%%%%%%%%%%%%%%%%%%%%%%%%%%%%%%%%%%%%%%%%%%%%%%%%%%%%%%%%%%%%%%%%%%%%%

\begin{document}

\maketitle

A todos los que de alguna manera estamos vinculados a la investigaci\'on en
matem\'aticas nos han preguntado alguna vez qu\'e se puede investigar
a\'un en matem\'aticas. La mayor\'{\i}a de la gente con conocimientos
matem\'aticos medios no piensa que en matem\'aticas quedan cosas por
descubrir. La manera m\'as f\'acil de responder a esa pregunta es mostrar
un problema con soluci\'on desconocida. Pero, claro, tiene que ser un problema
con un planteamiento lo suficientemente sencillo para que pueda ser entendido
por un profano. Hasta hace a\~nos, el ejemplo m\'as t\'{\i}pico era el
\'ultimo teorema de Fermat. Como ya est\'a demostrado, ha dejado de sernos
\'util en este sentido.

Afortunadamente, la teor\'{\i}a de n\'umeros tiene much\'{\i}simos otros
problemas abiertos de enunciado sencillo. Dos de los m\'as conocidos son la
conjetura de Goldbach (cada n\'umero par se puede expresar como suma de dos
primos) y la existencia de infinitos primos gemelos (primos que se diferencian
en dos). Realmente, los hay a cientos; tres libros dedicados a mostrar
problemas de este tipo son~\cite{Sier,Kl-Wa,Guy}. Es de destacar,
adem\'as, que muchos de ellos est\'an bien adaptados a la
realizaci\'on de experimentos con ordenadores. Si una conjetura dice que no
existen n\'umeros con ciertas propiedades y logramos encontrar uno, esto,
obviamente, echa por tierra la conjetura. En cambio, si la
suposici\'on es que existen infinitos n\'umeros que cumplen algo, con un
ordenador se puede intentar encontrar muchos de ellos, estudiando su
distribuci\'on, o buscando cada vez el mayor n\'umero que satisface la
propiedad. Estos r\'ecords, aunque no suelen aparecer en los medios de
comunicaci\'on, tienen exactamente la misma utilidad que los r\'ecords
deportivos, a los que tanto dinero se destina.

Aqu\'{\i} vamos a dedicarnos a hablar de uno de estos problemas abiertos, el
de las sucesiones alicuatorias. En el camino, nos aparecer\'an unos cuantos
m\'as.

Para un entero positivo~$n$ denotamos por $\sigma(n)$ la suma de todos sus
divisores (incluido $1$ y $n$). Si la descomposici\'on de~$n$ en factores
primos es $n = p_1^{a_1} \cdots p_k^{a_k}$, entonces
\begin{equation}
\label{ec:sigmaSimple}
    \sigma(p_1^{a_1} \cdots p_k^{a_k})
    = (1+p_1+\cdots+p_1^{a_1}) \cdots (1+p_k+\cdots+p_k^{a_k}).
\end{equation}
Esto es as\'{\i} porque, si desarrollamos la expresi\'on de la derecha,
aparecen como sumandos todos los divisores de~$p_1^{a_1} \cdots p_k^{a_k}$.
Adem\'as, $1+p+\cdots+p^a = \frac{p^{a+1}-1}{p-1}$
luego tenemos la siguiente f\'ormula directa para calcular $\sigma(n)$ en
funci\'on de la descomposici\'on en factores de~$n$:
\begin{equation}
\label{ec:sigma}
    \sigma(n) = \frac{p_1^{a_1+1}-1}{p_1-1} \cdots
      \frac{p_k^{a_k+1}-1}{p_k-1}.
\end{equation}

A partir de~$\sigma$ construimos la funci\'on $s(n) = \sigma(n) - n$.
Obviamente, $s(n)$ es la suma de los divisores propios de~$n$ (esto es,
excluyendo el mismo~$n$). Con otras palabras, $s(n)$ es la suma de
las partes al\'{\i}cuotas de~$n$.

Relacionados con~$s$ (o con~$\sigma$) se pueden definir varios tipos de
n\'u\-me\-ros (un libro con resultados recientes sobre el tema, y que
incide en su aspecto computacional, es~\cite{Yan}):

%\begin{itemize}

%\item
\textbf{N\'umeros perfectos}: Un n\'umero se dice perfecto si es igual a
la suma de sus divisores propios, es decir, si $s(n) = n$ (o $\sigma(n) =
2 \cdot n$). Si $s(n) > n$, el n\'umero $n$ se llama abundante;
si $s(n) < n$, defectuoso. Los primeros n\'umeros perfectos son $6$, $28$ y
$496$. Es bien conocido que $n$ es un n\'umero perfecto par si y s\'olo si
es de la forma $n = 2^{p-1}(2^p-1)$ con $2^p-1$ primo (lo cual, a su vez,
requiere que~$p$ sea primo). Pero se desconoce si existen o no n\'umeros
perfectos impares. En cuanto a los pares, encontrarlos se reduce a la
b\'usqueda de n\'umeros $M_p = 2^p-1$ (que se denominan n\'umeros de
Mersenne) primos. No se sabe si hay o no infinitos de ellos; ni tampoco si hay
infinitos compuestos. Actualmente, se conocen~$37$; el mayor es $2^{3021377}-1$,
que tiene m\'as de $900.000$ cifras (existe un premio de $50.000$ d\'olares para
la primera persona que encuentre un primo de un mill\'on de cifras). La
b\'usqueda de n\'umeros de Mersenne primos est\'a coordinada por G.~Woltmann con
el proyecto \textit{Great Internet Mersenne Prime Search} en
\url{http://www.mersenne.org/prime.htm}. All\'{\i} se pueden conocer los
\'ultimos avances e, incluso, descargar los programas necesarios para contribuir
en la b\'usqueda con nuestro ordenador personal. Los tres \'ultimos primos de
Mersenne conocidos han sido encontrados por colaboradores de este proyecto en
ordenadores personales (el anterior hab\'{\i}a sido encontrado en un
superordenador CRAY).

%\item
\textbf{N\'umeros amigos}: Son parejas de n\'umeros $(n,m)$ tales que 
$s(n)=m$ y $s(m)=n$ (o, lo que es equivalente, $\sigma(n) = \sigma(m) = n+m$).
Ejemplos de amigos son $(220,284)$, $(2620,2924)$ y $(5020,5564)$.
Hasta ahora, se conocen varios miles. Se han desarrollado diversos m\'etodos
que, a partir de una pareja de amigos, generan otras parejas de n\'umeros que
son amigos con bastante probabilidad. Cuando se han aplicado, mediante el uso
de ordenador, a tablas de amigos ya conocidos, se han obtenido m\'as nuevos
que los ya conocidos (v\'ease, por ejemplo~\cite{teRie}). Esto es un argumento
heur\'{\i}stico en favor de la existencia de infinitos pares de amigos.
Pero tampoco se tiene una demostraci\'on.

%\item
\textbf{Ciclos o n\'umeros pandilla}: Son tuplas $(a_1, \dots, a_l)$
tales que $s(a_1)=a_2$, $s(a_2)=a_3$, \dots, $s(a_l)=a_1$. 
Hasta ahora, se han encontrado~$41$ ciclos de longitud~$l$ mayor
que~$2$ (obviamente, los de longitud uno son los n\'umeros perfectos y los
de longitud dos son los amigos). S\'olo se conocen ciclos de longitudes $4$,
$5$, $6$, $8$, $9$ y $28$. La mayor\'{\i}a de ellos han sido hallados con
ordenador. Curiosamente no sucedi\'o as\'{\i} con el que, a priori,
podr\'{\i}a parecer m\'as raro, el de longitud $28$: fue encontrado por
Poulet en $1918$; su componente m\'as peque\~no es $14316$. Se desconoce,
por ejemplo, si existen ciclos de longitud~$3$. O si existen infinitos.

%\item
\textbf{Intocables de Erd\H{o}s}: Se llaman as\'{\i} a los~$n$ tales que
$n \ne s(m)$ para todo~$m$. Por ejemplo, $2$ y~$5$. De nuevo, no se sabe si
hay o no infinitos de ellos.%
\footnote{Nota posterior: En~\cite[problema 91, p.~119]{Sier}, la existencia o no
de infinitos n\'umeros intocables aparece como problema abierto. Pero, en
1973, el mismo Erd\H{o}s prob\'o no s\'olo que existen infinitos de ellos,
sino algo m\'as potente; concretamente, que su densidad inferior es positiva.
V\'ease P.~Erd\H{o}s,
\textit{\"Uber die Zhalen der Form $\sigma(n)-n$ und $n-\varphi(n)$},
\textrm{Elem.\ Math.}
\textbf{28} (1973), 83--86; MR \textbf{49} \#2502.}

%\end{itemize}

%%%%%%%%%%%%%%%%%%%%%%%%%%%%%%%%%%%%%%%%%%%%%%%%%%%%%%%%%%%%%%%%%%%%%%%%%%%
\section*{Sucesiones alicuatorias}
%%%%%%%%%%%%%%%%%%%%%%%%%%%%%%%%%%%%%%%%%%%%%%%%%%%%%%%%%%%%%%%%%%%%%%%%%%%

Ya estamos en condiciones de poder definir las \textsl{sucesiones de sumas de
partes al\'{\i}cuotas} o, para abreviar,
\textsl{sucesiones alicuatorias}:
Dado un entero positivo~$n$ construimos la sucesi\'on
$\{s^k(n)\}_{k \in \mathbb{N}}$ definida recursivamente mediante
$s^1(n) = s(n)$, $s^2(n) = s(s^1(n))$, \dots, $s^k(n) = s(s^{k-1}(n))$.

Para una de tales sucesiones, existen cuatro posibilidades: (a)~Que la
sucesi\'on termine en~$1$ (siendo el n\'umero anterior un primo). (b)~Que la
sucesi\'on llegue a un n\'umero perfecto (y a partir de entonces permanezca
constante). (c)~Que llegue a un par de amigos o un ciclo. (d)~Que no est\'e
acotada.

E.~Catalan~\cite{Cat} en $1887$ y L.~E.~Dickson~\cite{Dick} en $1913$
conjeturaron que la po\-si\-bi\-li\-dad~(d) nunca ocurr\'{\i}a. La conjetura de
Catalan-Dickson contin\'ua sin ser demostrada o refutada.

Una manera de ver que es falsa ser\'{\i}a encontrar un~$n$ tal que $s^k(n) >
s^{k-1}(n)$ para todo~$k$. Aunque parece dif\'{\i}cil que tal~$n$ exista.
S\'{\i} que existen sucesiones con tantos t\'erminos crecientes como
se desee, tal como prob\'o H. W. Lenstra (v\'ease~\cite{Erdos}):
para cada~$k$, existe~$n$ tal que $n < s(n) < \cdots < s^k(n)$. El m\'etodo
es constructivo y f\'acil de entender, aunque los n\'umeros que van saliendo
son enormes. Argumentos heur\'{\i}sticos de Guy-Selfridge ---basados en la
persistencia de ciertos patrones que se repiten en las sucesiones--- hacen suponer
que, para infinitos valores de~$n$, esto se puede conseguir con $k < (\log
n)^{1-\varepsilon}$. Es m\'as, ellos se atreven a conjeturar en el sentido
contrario al de Catalan-Dickson: las sucesiones $\{s^k(n)\}_{k \in \mathbb{N}}$
son no acotadas para casi todo~$n$ par; es decir, la proporci\'on de los enteros
pares para los cuales $\{s^k(n)\}_{k \in \mathbb{N}}$ est\'a acotada, tiende a
cero.

Antes de continuar, veamos un ejemplo de este tipo de patrones que se van
repitiendo a lo largo de los t\'erminos de la sucesi\'on; al menos, que se repiten
con bastante probabilidad.

Supongamos $n = 2^3 \cdot 3 \cdot 5 \cdot p \cdot q$, con~$p$ y~$q$ primos
distintos mayores que~$5$; lo mismo ocurre si hay m\'as factores primos pero, por
simplicidad, ve\'amoslo de esa manera. Por~(\ref{ec:sigmaSimple}),
$s(n) = (1+2+4+8) \cdot 4 \cdot 6 \cdot (p+1) \cdot (q+1) - 2^3 \cdot 3 \cdot 5
\cdot p \cdot q = 15 \cdot 2^2 \cdot 2 \cdot 3 \cdot (p+1) \cdot (q+1)
- 2^3 \cdot 3 \cdot 5 \cdot p \cdot q = 2^3 \cdot 3 \cdot 5 \cdot
[3 \cdot (p+1) \cdot (q+1) - p \cdot q] = 2^3 \cdot 3 \cdot 5 \cdot m$,
con~$m$ impar y no m\'ultiplo de~$3$. Es decir, de nuevo obtenemos el patr\'on
$2^3 \cdot 3 \cdot 5$, con~$2$ y~$3$ elevados exactamente a la misma potencia. Sin
embargo, no se garantiza que~$5$ vuelva a aparecer como factor de $s(n)$
elevado \'unicamente a la potencia~$1$. El comportamiento es similar si los primos
mayores que~$5$ est\'an elevados a potencias impares, ya que $1 + p + \cdots +
p^a$ es par cuando~$a$ es impar. En cambio, un factor $p^a$ con~$a$ par no
contribuye con un factor~$2$ en el primer sumando de la expresi\'on entre
corchetes que da~$m$, puesto que $1 + p + \cdots + p^a$ es impar; es como si~$n$
no tuviese el factor~$p^a$.

Es importante se\~nalar que este tipo de patrones tambi\'en pueden desaparecer.
Por ejemplo, siguiendo con el ejemplo anterior, as\'{\i} puede suceder si~$n$
tiene la forma $n = 2^3 \cdot 3 \cdot 5^2 \cdot p$ (un solo primo~$p$ mayor
que~$5$). Pero esta combinaci\'on de factores no es muy probable.

Las estructuras estables no son raras. Ocurren siempre que aparece como
factor un n\'umero perfecto (por ejemplo, $n = 2 \cdot 3 \cdot p \cdot q$ o
$n = 2^2 \cdot 7 \cdot p \cdot q$) y con algunas otras combinaciones de
primos, como $2^3 \cdot 3$ o $2^3 \cdot 3 \cdot 5$. Las comprobaciones son
similares. No incidiremos m\'as en ellas. Simplemente, nos limitamos a comentar
que su existencia ha sido analizada rigurosamente y plasmada en forma de
definiciones generales. As\'{\i}, Guy y Selfridge
\cite{Guy-Sel} dan los conceptos de gu\'{\i}a y conductor (en ingl\'es,
\textit{guide} y \textit{driver}).

%\section*{Uno poco de historia}

El primer n\'umero que ofreci\'o serias dudas en cuanto al comportamiento de la
sucesi\'on alicuatoria que comienza en \'el fue $n = 138$ (Poulet, $1918$). Fue
resuelto por D. H. Lehmer, que encontr\'o $s^{177}(138) = 1$, pasando por un
m\'aximo $s^{117}(138)$ de $12$ cifras (en notaci\'on decimal). Desde entonces, el
menor~$n$ cuya sucesi\'on tiene comportamiento incierto es $276$. Del estudio
experimental de tal sucesi\'on, computando cada vez m\'as t\'erminos, se han
ocupado G.~A.~Paxson, H.~Cohen, D.~H.~Lehmer, H.~J.~Godwin, J.~L.~Selfridge,
M.~C.~Wunderlich, T.~Struppeck, R.~K.~Guy, A.~Guy, M.~Dickerman, P.~Zimmermann y
los autores. Actualmente, se ha llegado hasta $s^{1119}(276)$, un n\'umero de
$107$ cifras (P.~Zimmermann, con la colaboraci\'on de S.~Wagstaff y A.~Lenstra).

Las sucesiones alicuatorias $\{s^k(n)\}_{k \in \mathbb{N}}$ y
$\{s^k(m)\}_{k \in \mathbb{N}}$ correspondientes a dos enteros positivos~$n$ y
$m$, con $n < m$, se dice que son laterales si existen \'{\i}ndices $i$ y~$j$
tales que $s^i(n) = s^j(m)$ (l\'ogicamente, a partir de entonces ambas
sucesiones coinciden); y se llama sucesi\'on principal a la correspondiente
a~$n$. Es claro que, en el estudio de las sucesiones alicuatorias, basta con
analizar las sucesiones principales.

En principio, el tratamiento computacional de una sucesi\'on
$\{s^k(n)\}_{k \in \mathbb{N}}$ para un~$n$ concreto parece f\'acil: basta
aplicar reiteradamente la f\'or\-mu\-la~(\ref{ec:sigma}). Pero la dificultad
radica en que la sucesi\'on puede alcanzar t\'erminos muy grandes. Para poder
seguir,~(\ref{ec:sigma}) requiere factorizar enteros enormes. Y, con
cualquiera de los algoritmos de factorizaci\'on conocidos, esto supone un
tiempo de c\'alculo que crece exponencialmente con el n\'umero de
d\'{\i}gitos del n\'umero a factorizar. 

Adem\'as de a la correspondiente a~$276$, se ha dedicado un gran esfuerzo
computacional a muchas otras sucesiones. Comentemos a continuaci\'on en qu\'e
se ha centrado fundamentalmente este trabajo computacional.
En~\cite[Cap.~B6]{Guy}, se pueden encontrar m\'as datos hist\'oricos
(hasta $1994$), y una amplia bibliograf\'{\i}a sobre el tema.

Las cinco de Lehmer: Son las sucesiones principales que empiezan por un~$n$
menor que $1000$. Concretamente, $276$, $552$, $564$, $660$ y~$966$.

Las catorce de Godwin: Son las que comienzan con~$n$ entre $1000$ y $2000$.
Actualmente, s\'olo $12$ de ellas permanecen en duda. Godwin, en $1980$,
encontr\'o el final de la sucesi\'on correspondiente a $1984$; result\'o ser
una sucesi\'on con un m\'aximo de~$29$ cifras y una longitud de $672$.
Y Dickermann, en $1994$, el final de la que comienza en $1248$;
esta vez, se obtuvo un m\'aximo de~$58$ cifras y una longitud de~$1075$.

Asimismo, los autores han estudiado todas las sucesiones que comienzan en
un n\'umero~$n$ menor de $10000$. En todas ellas, se ha avanzado hasta
t\'erminos mayores de $10^{90}$ y se ha encontrado el final de algunas que
anteriormente estaban en duda. M\'as adelante en este mismo art\'{\i}culo
comentamos con mayor detalle nuestro trabajo con estas sucesiones.

Por \'ultimo, W.~Creyaufm\"uller~\cite{Crey} est\'a intentando clasificar todas
las sucesiones que comienzan en un $n < 10^6$. El estudio de las sucesiones
que ofrecen dudas se prosigue hasta alcanzar t\'erminos de~$60$
cifras. \'Este es un proyecto que est\'a a\'un en desarrollo.
Puede encontrarse informaci\'on actualizada en
\url{http://www.loria.fr/~zimmerma/records/aliquot.html}
(p\'agina mantenida por P.~Zimmermann).

Adem\'as de esta p\'agina web, existen otras dos dedicadas a los progresos en
sucesiones alicuatorias: La mantenida por uno de los autores,
\url{http://www.unirioja.es/dptos/dmc/jvarona/aliquot.html},
y la de W. Bosma,
\url{http://www-math.sci.kun.nl/math/~bosma/nuth/ali.html}.

Todo el esfuerzo computacional para ir calculando t\'erminos de una
sucesi\'on alicuatoria se hace aplicando las
f\'ormulas~(\ref{ec:sigmaSimple}) o~(\ref{ec:sigma}). Esto, cuando la
sucesi\'on alcanza t\'erminos enormes, requiere un gran tiempo de
factorizaci\'on. Las mejoras en los algoritmos de factorizaci\'on (y la
velocidad de los ordenadores) permiten avanzar en el estudio de las
sucesiones.

Pero nadie puede garantizar que no exista alg\'un m\'etodo r\'apido que
permita calcular $s(n)$ (o $\sigma(n)$) a partir de~$n$ sin necesidad de
factorizar~$n$. Realmente, ya Euler encontr\'o un algoritmo recursivo para el
c\'alculo de~$\sigma(n)$. Lamentablemente, la recursi\'on, cuando~$n$ es
grande, requiere una gran cantidad de pasos previos. El m\'etodo puede ser
descrito mediante
\begin{equation*}
\begin{align*}
    \sigma(n) &= \sigma(n-1) + \sigma(n-2) - \sigma(n-5) - \sigma(n-7) \\
      &+ \sigma(n-12) + \sigma(n-15) - \sigma(n-22) - \sigma(n-26) \\
      &+ \sigma(n-35) + \sigma(n-40) - \sigma(n-51) - \sigma(n-57) \\
      &+ \sigma(n-70) + \sigma(n-77) - \sigma(n-92) - \sigma(n-100) \\
      &+ \cdots .
\end{align*}
\end{equation*}
No haremos m\'as hincapi\'e en esto. Para conocer m\'as detalles sobre
esta f\'ormula, cons\'ultese~\cite[Cap.~6]{Pol}.

%%%%%%%%%%%%%%%%%%%%%%%%%%%%%%%%%%%%%%%%%%%%%%%%%%%%%%%%%%%%%%%%%%%%%%%%%%%
\section*{Nuestros progresos}
%%%%%%%%%%%%%%%%%%%%%%%%%%%%%%%%%%%%%%%%%%%%%%%%%%%%%%%%%%%%%%%%%%%%%%%%%%%

Tal como coment\'abamos anteriormente, los autores se han preocupado en
estudiar computacionalmente las sucesiones alicuatorias que comienzan en
$n < 10000$. Este trabajo avanzaba no s\'olo en el estudio de las sucesiones
de Lehmer y Godwin, sino tambi\'en en estudios previos de otros autores, como
A.~Guy y R.~K.~Guy, que hab\'{\i}an analizado las sucesiones que comenzaban en
$n \le 7044$. Fruto de este trabajo, hallamos el final de varias sucesiones.
En particular, el de la sucesi\'on que comienza en $4170$. \'Esta es, hasta
ahora, la sucesi\'on con final conocido que alcanza un t\'ermino mayor;
concretamente $s^{289}(4170)$ tiene~$84$ cifras y acaba en
$s^{869}(4170) = 1$, siendo~$79$ el primo previo. (Obviamente, esta
afirmaci\'on requiere descartar casos triviales: por ejemplo, cualquier
sucesi\'on que comience en un primo acaba inmediatamente por grande que sea el
primo.) Este trabajo fue publicado en~\cite{Ben-Var}. All\'{\i}, adem\'as,
se muestra el estado de todas las sucesiones principales que comienzan en
$n < 10000$ hasta alcanzar t\'erminos de, al menos,~$75$ d\'{\i}gitos.

\begin{table}
%\renewcommand{\topfraction}{0.95}
%\renewcommand{\bottomfraction}{0.95}
%\renewcommand{\textfraction}{0.05}
\begin{center}
\renewcommand{\arraystretch}{0.80}
%\small % To choose: \small, \footnotesize or \scriptsize.
%\begin{tabular}{|c|c@{\,}c@{\,}c@{\,}c@{\,}c@{\,}c@{\,}c@{\,}c@{\,}c|}
\begin{tabular}{|c|ccccccccc|}
\hline
\info &
\num{276}{1119}{107}{2 \pp 3} &
\num{552}{808}{114}{2^5 \pp 3 \pp 7} &
\num{564}{2970}{102}{2^2 \pp 7} &
\num{660}{439}{110 \ldw{104}}{2^2 \nodriver} & % was 2^5*3*7 at i371
\num{966}{472}{108}{2 \pp 3} & %%% 2 \pp 3^2
\num{1074}{1543}{103}{2^2 \pp 7} &
\num{1134}{2234}{121}{2^5 \pp 3 \pp 7} & %%% 2^5 \pp 3^3 \pp 7
\num{1464}{1815}{97}{2^2 \pp 7} &
\num{1476}{1055}{106}{2^3 \pp 3 \pp 5}
\\ \hline
\info &
\num{1488}{748}{97 \ldw{94}}{2^3 \pp 5 \nodriver} &
\num{1512}{1602}{94}{2^2 \pp 7} &
\num{1560}{1336}{101}{2^5 \pp 3 \pp 7} & %%% 2^5 \pp 3^2 \pp 7
\num{1578}{1084}{100 \ldw{88}}{2^4 \nodriver} &
\num{1632}{713}{102}{2^3 \pp 3 \pp 5} & %%% 2^3 \pp 3^2 \pp 5
\num{1734}{1367}{99}{2^2 \pp 7} &
\num{1920}{1931}{97}{2^2 \pp 7} &
\num{1992}{985}{102}{2^3 \pp 3 \pp 5} & %%% 2^3 \pp 3^3 \pp 5
\num{2232}{390}{102 \ldw{95}}{2^6 \nodriver}
\\ \hline
\info &
\num{2340}{471}{99}{2^3 \pp 3 \pp 5} &
\num{2360}{974}{95}{2^2 \pp 7} &
\num{2484}{796}{97}{2^3 \pp 3} & %%% 2^3 \pp 3^3
\num{2514}{2836}{95}{2^3 \pp 3 \pp 5} & %%% 2^3 \pp 3 \pp 5^2
\num{2664}{761}{100}{2^2 \pp 7} &
\num{2712}{1347}{95}{2^2 \pp 7} &
\num{2982}{810}{97}{2^4 \pp 31} &
\num{3270}{417}{98}{2^5 \pp 3 \pp 7} & %%% 2^5 \pp 3 \pp 7^2
\num{3366}{1062}{100 \ldw{82}}{2^3 \nodriver}
\\ \hline
\info &
\num{3408}{840}{95}{2^3 \pp 3 \pp 5} &
\num{3432}{933}{103}{2^3 \pp 3 \pp 5} &
\num{3564}{779}{100}{2^3 \pp 3} & %%% 2^3 \pp 3^2
\num{3630}{1193}{94 \ldw{94}}{2^4 \nodriver} &
\num{3678}{1201}{98}{2^2 \pp 7} &
\num{3774}{1193}{98 \ldw{89}}{2^7 \pp 3 \nodriver} &
\num{3876}{830}{96}{2^2 \pp 7} &
\num{3906}{679}{94}{2^2 \pp 7} &
\num{4116}{897}{97 \ldw{92}}{2^3 \nodriver}
\\ \hline
\info &
\num{4224}{519}{98}{2^3 \pp 3 \pp 5} &
\num{4290}{913}{92}{2^2 \pp 7} &
\num{4350}{1165}{97}{2^2 \pp 7} &
\num{4380}{965}{100}{2^2 \pp 7} & %%% 2^2 \pp 7^2
\num{4788}{2152}{105}{2^3 \pp 3 \pp 5} & %%% 2^3 \pp 3^3 \pp 5
\num{4800}{1135}{101}{2^4 \pp 31} &
\num{4842}{440}{91}{2^2 \pp 7} &
\num{5148}{1545}{95}{2 \pp 3} &
\num{5208}{1710}{96}{2^3 \pp 3 \pp 5}
\\ \hline
\info &
\num{5250}{1490}{93 \ldw{89}}{2^5 \nodriver} &
\num{5352}{683}{93}{2^2 \pp 7} &
\num{5400}{2696}{93}{2^2 \pp 7} &
\num{5448}{1185}{96}{2^3 \pp 3 \pp 5} &
\num{5736}{1093}{100}{2^2 \pp 7} &
\num{5748}{958}{91}{2^2 \pp 7} &
\num{5778}{742}{95}{2^4 \pp 31} &
\num{6160}{1630}{96 \ldw{84}}{2^2 \nodriver} &
\num{6396}{1234}{92}{2^3 \pp 3 \pp 5} %%% 2^3 \pp 3 \pp 5^2
\\ \hline
\info &
\num{6552}{893}{93}{2^3 \pp 3} &
\num{6680}{1815}{94}{2^4 \pp 31} &
\num{6822}{1177}{97}{2^4 \pp 31} &
\num{6832}{885}{104}{2^3 \pp 3} &
\num{6984}{1764}{96}{2^4 \pp 31} &
\num{7044}{1113}{102}{2^4 \pp 31} &
\num{7392}{498}{96}{2^3 \pp 3 \pp 5} & %%% 2^3 \pp 3^2 \pp 5^2
\num{7560}{846}{97 \ldw{86}}{2^3 \pp 5 \nodriver} &
\num{7890}{891}{99}{2^2 \pp 7} %%% 2^2 \pp 7^2
\\ \hline
\info &
\num{7920}{951}{95}{2^6 \pp 127} &
\num{8040}{2205}{94}{2^3 \pp 3 \pp 5} &
\num{8154}{647}{96}{2^2 \pp 7} &
\num{8184}{1241}{102}{2^2 \pp 7} &
\num{8288}{823}{98 \ldw{95}}{2^2 \nodriver} &
\num{8352}{1246}{93}{2^2 \pp 7} &
\num{8760}{2145}{94}{2^3 \pp 3 \pp 5} & %%% 2^3 \pp 3^3 \pp 5^3
\num{8844}{1184}{101}{2^4 \pp 31} &
\num{8904}{963}{95}{2^2 \pp 7}
\\ \hline
\info &
\num{9120}{532}{92}{2^3 \pp 3} &
\num{9282}{505}{95}{2^3 \pp 3} &
\num{9336}{608}{97}{2^6 \pp 127} &
\num{9378}{2111}{92}{2^2 \pp 7} &
\num{9436}{545}{94}{2 \pp 3} &
\num{9462}{447}{97}{2^3 \pp 3 \pp 5} &
\num{9480}{924}{91}{2 \pp 3} &
\num{9588}{1492}{93 \ldw{90}}{2^3 \nodriver} &
\num{9684}{617}{92}{2^5 \pp 3 \pp 7}
\\ \hline
\info &
\num{9708}{671}{95}{2^2 \pp 7} &
\num{9852}{591}{92}{2^2 \pp 7} &&&&&&&
\\ \hline
\end{tabular}
\end{center}
\caption{Sucesiones alicuatorias cuyo final est\'a en duda.} \label{tabla}
\end{table}
%%%%%%%%%%%%%%%%%%%%%%%%%%%%%%%%%%%%%%

No entraremos en detalles de los algoritmos de factorizaci\'on empleados.
Simplemente, comentar que, principalmente, se ha usado una combinaci\'on del
m\'etodo de las curvas el\'{\i}pticas y de la criba cuadr\'atica
multipolinomial (v\'ease~\cite{Cohen}).
As\'{\i}, para descomponer en factores un n\'umero de~$60$ cifras (sin
factores peque\~nos) se emplea alrededor de tres cuartos de hora en un
Pentium~100; de modo similar, hora y media si el n\'umero tiene~$65$ cifras,
$5$~horas para uno con $70$ cifras, $18$ horas con $75$ cifras, un d\'{\i}a y
medio con $80$ cifras, $4$~d\'{\i}as con $85$ cifras, y dos semanas con $90$
cifras.
Aumentar el n\'umero de cifras del \'ultimo t\'ermino alcanzado en una
sucesi\'on alicuatoria puede requerir muchas iteraciones, y estamos tratando m\'as
de~$80$ sucesiones. Esto supone una gran cantidad de tiempo de c\'alculo.

Es de destacar que en ning\'un momento hemos utilizado ning\'un programa comercial
para llevar a cabo los c\'omputos (la mayor\'{\i}a de los paquetes inform\'aticos
comerciales destinados al c\'alculo matem\'atico, tanto num\'erico como
simb\'olico, ni siquiera tienen los algoritmos necesarios implementados y, si los
tienen y hemos conseguido probarlos, sus implementaciones han resultado ser
bastante peores que los de diversos programas que se pueden encontrar gratis en
internet). A lo largo de diversas etapas de nuestro trabajo, los paquetes que
hemos empleado han sido
PARI\footnote{C.~Batut, D.~Bernardi, H.~Cohen y M.~Olivier, disponible en
\url{ftp://megrez.ceremab.u-bordeaux.fr/pub/pari}.},
KASH\footnote{Desarrollado por M.~E.~Posh y \textit{The KANT Group}; 
disponible en \url{ftp://ftp.math.tu-berlin.de/pub/algebra/Kant/Kash}.
Para m\'as informaci\'on, v\'ease
\url{http://www.math.TU-Berlin.DE/~kant/kash.html}.},
MIRACL\footnote{M.~Scott, disponible en
\url{ftp://ftp.compapp.dcu.ie/pub/crypto}.} y, fundamentalmente,
UBASIC\footnote{Yuji Kida, disponible en \url{ftp://rkmath.rikkyo.ac.jp/ubibm}.}.
Los c\'alculos se han llevado a cabo en tiempos muertos, noches y fines de semana
de numerosos ordenadores de las instituciones a las que pertenecemos, y en los
personales de los autores o sus amigos.

M\'as tarde, tuvimos conocimiento de que P.~Zimmermann y W.~Bosma estaban
haciendo, en parte, el mismo trabajo. Concretamente, P.~Zimmermann estaba
analizando las cinco sucesiones de Lehmer y las que comienzan por $1074$ y
$1134$, y W.~Bosma se estaba dedicando a las sucesiones que comienzan en $n \le
50000$ (por ejemplo, el final de la sucesi\'on $3556$ fue encontrado por Bosma y
los autores independientemente); l\'ogicamente, cuantas m\'as sucesiones se
abarcan, menos tiempo computacional es posible dedicar a cada una. Entonces, nos
repartimos el trabajo. Concretamente, en lo referente a las sucesiones
que comienzan por $n < 10000$, en las que nos centraremos aqu\'{\i}, P.~Zimmermann
sigui\'o con las suyas y nosotros nos dedicamos al resto. Realmente, desde la
publicaci\'on de~\cite{Ben-Var} no se ha encontrado el final de ninguna nueva de
entre ellas. Pero se ha avanzado bastante en todas.

El estado actual de los c\'alculos para sucesiones alicuatorias que comienzan en
$n < 10000$ es el que aparece resumido en la tabla~\ref{tabla}, que ampl\'{\i}a
los resultados de~\cite{Ben-Var}. En ella mostramos las~$83$ sucesiones
principales cuyo final es a\'un desconocido. Tambi\'en incluimos el n\'umero de
d\'{\i}gitos decimales del \'ultimo $s^k(n)$ alcanzado para cada sucesi\'on.
Los datos de las~$7$ primeras sucesiones son de
P.~Zimmermann. Es de destacar que en todas las sucesiones se ha llegado a
t\'erminos mayores que $10^{90}$; en varias de ellas se han alcanzado las~$100$
cifras. El desarrollo completo de todas las sucesiones que parecen en la tabla
puede conseguirse, mediante \texttt{ftp} an\'onimo, en
\url{ftp://navalsaz.unirioja.es/pub/aliquot}.

%%%%%%%%%%%%%%%%%%%%%%%%%%%%%%%%%%%%%%%%%%%%%%%%%%%%%%%%%%%%%%%%%%%%%%%%%%%
\subsection*{Agradecimientos}
%%%%%%%%%%%%%%%%%%%%%%%%%%%%%%%%%%%%%%%%%%%%%%%%%%%%%%%%%%%%%%%%%%%%%%%%%%%
Deseamos dar las gracias a Paul Zimmermann por brindarnos su colaboraci\'on
durante varios a\~nos, y por permitirnos reproducir aqu\'{\i} los \'ultimos
t\'erminos alcanzados por \'el en las sucesiones que comienzan en $n \le 1134$
(que tambi\'en se pueden encontrar en su p\'agina web antes citada).
As\'{\i} mismo, debemos nuestro agradecimiento a las diversas personas
que nos han cedido tiempo de c\'alculo en sus ordenadores personales o en los
de sus instituciones; entre ellos Carlos Jos\'e P\'erez, Emilio
Fern\'andez, Jes\'us Murillo, M.\textsuperscript{a} Vico Pascual, Julio Rubio,
Lourdes Benito y Esperanza Infante. Tambi\'en a numerosos becarios de salas
inform\'aticas de la Universidad de La Rioja. Si alguien se siente olvidado,
le pedimos disculpas de antemano.

Del mismo modo, gracias a la Universidad de La Rioja y al Instituto Sagasta de
Logro\~no por permitirnos usar sus ordenadores, en los que se llevaron a cabo
la mayor parte de los c\'alculos. Por \'ultimo, el segundo autor quiere
agradecer a la Universidad de La Rioja los innumerables cortes de electricidad
que hicieron el trabajo computacional mucho m\'as entretenido.

%%%%%%%%%%%%%%%%%%%%%%%%%%%%%%%%%%%%%%%%%%%%%%%%%%%%%%%%%%%%%%%%%%%%%%%%%%%
\subsection*{Despedida}
%%%%%%%%%%%%%%%%%%%%%%%%%%%%%%%%%%%%%%%%%%%%%%%%%%%%%%%%%%%%%%%%%%%%%%%%%%%
Hemos comenzado este art\'{\i}culo hablando de problemas abiertos en
teor\'{\i}a de n\'umeros. Acabemos con un acertijo relacionado con ellos.

Supongamos que, cada mes, los matem\'aticos plantean un nuevo problema de
teor\'{\i}a de n\'umeros. Y que, cada a\~no, se resuelve uno de los a\'un 
abiertos, siempre el m\'as antiguo. ?`Cu\'antos problemas quedar\'an sin
resolver al final de los tiempos?
?`Y si, en vez de resolverse uno cada a\~no, se resuelven dos, siempre los
dos que ocupan las posiciones centrales (considerando los problemas ordenados
temporalmente)?

%%%%%%%%%%%%%%%%%%%%%%%%%%%%%%%%%%%%%%%%%%%%%%%%%%%%%%%%%%%%%%%%%%%%%%%%%%%
\begin{thebibliography}{99}
%%%%%%%%%%%%%%%%%%%%%%%%%%%%%%%%%%%%%%%%%%%%%%%%%%%%%%%%%%%%%%%%%%%%%%%%%%%

\bibitem{Ben-Var} M. Benito y J. L. Varona,
Advances in aliquot sequences,
\textit{Math. Comp.}
\textbf{68} (1999), 389--393.

\bibitem{Cat} E. Catalan,
Propositions et questions diverses,
\textit{Bull. Soc. Math. France}
\textbf{18} (1887--88), 128--129.

\bibitem{Cohen} H. Cohen,
\textit{A course in computational algebraic number theory},
Springer-Verlag, 1993.

\bibitem{Crey} W. Creyaufm\"uller,
\textit{Primzahlfamilien} (2.\textsuperscript{a}~ed.),
Wolfgang Creyaufm{\"u}ller, Aachen, 1997.
% En alem\'an, 262 p\'aginas.

\bibitem{Dick} L. E. Dickson,
Theorems and tables on the sum of the divisors of a number,
\textit{Quart. J. Math.}
\textbf{44} (1913), 264--296.

\bibitem{Erdos} P. Erd\H{o}s,
On asymptotic properties of aliquot sequences,
\textit{Math. Comp.}
\textbf{30} (1976), 641--645.
%\MR{53:7919}

% \bibitem{Guy-Guy} A. W. P. Guy y R. K. Guy,
% A record aliquot sequence,
% en \textit{Computation 1943--1993: A Half-Century of
% Computational Mathematics} (Vancouver, 1993),
% \textit{Proc. Sympos. Appl. Math.}
% \textbf{48} (1994),
% Amer. Math. Soc., Providence, RI, 1994,
% pp.~557--559.
% %\MR{95k:11164}

\bibitem{Guy} R. K. Guy,
\textit{Unsolved Problems in Number Theory} (2.\textsuperscript{a}~ed.),
Springer-Verlag, 1994.
%\MR{96e:11002}

\bibitem{Guy-Sel} R. K. Guy y J. L. Selfridge,
What drives an aliquot sequence?,
\textit{Math. Comp.}
\textbf{29} (1975), 101--107.
Corrigendum, \textit{ibid.} \textbf{34} (1980), 319--321.
%\MR{52:5542}; \MR{81f:10008}

\bibitem{Kl-Wa} V. Klee y S. Wagon,
\textit{Old and new unsolved problems in plane geometry and number theory},
The Math. Assoc. of Amer., Washington, 1991.

\bibitem{Pol} G. P\'olya,
\textit{Matem\'aticas y razonamiento plausible},
Tecnos, Madrid, 1966.

% \bibitem{Riesel} H. Riesel,
% \textit{Prime numbers and computer methods for factorization},
% Birkh\"auser, 1987. %% NO TIENE EL MPQS

\bibitem{teRie} H. J. J. te Riele,
On generating new amicable pairs from given amicable pairs,
\textit{Math. Comp.}
\textbf{42} (1984), 219--223.

\bibitem{Sier} W. Sierpi\'nski,
\textit{A selection of problems in the theory of numbers},
Pergamon Press, 1964.

\bibitem{Yan} S. Y. Yan,
\textit{Perfect, amicable and sociable numbers. A computacional approach},
World Scientific, Singapur, 1996.

\end{thebibliography}

%%%% DIRECCION DE LOS AUTORES DEL ART\'ICULO

\vskip 1pt
\begin{flushright}
{\small Manuel Benito,\\
Instituto Sagasta,\\
C/ Doctor Zub\'{\i}a s/n, 26003 Logro\~no.\\
e-mail: \url{mbenit8@palmera.pntic.mec.es}}
\end{flushright}

\begin{flushright}
{\small Juan Luis Varona,\\
Departamento de Matem\'aticas y Computaci\'on,
Universidad de La Rioja,\\
C/ Luis de Ulloa s/n, 26004 Logro\~no.\\
e-mail: \url{jvarona@dmc.unirioja.es}\\
URL: \url{http://www.unirioja.es/dptos/dmc/jvarona/welcome.html}}
\end{flushright}

\end{document}

 